\section{Coordinate variation and final-state marginals}
\label{sec:peps_regularized_patch_stationarity}

Retain the regions, identity extensions, and ordered product of
Definition~\ref{def:peps_regularized_patch_objective}.
Write $K_j=(x_j+bI_{X_j})^{-a_j/2}$, so $L_j=\widehat{K_j}_{X_j}$.
The coordinate variation below is the unstripped first variation in the
\href{https://github.com/openai/math/blob/adc7f1241b42e322a6451854ab7e4b4c146bf78a/preprints/Polynomial-PEPS-approximation-of-gapped-square-grid-ground-states-September-24-2026/build/sections/03-patches.tex\#L134-L149}{proof of the regularized patch estimate}.
Its marginals are those of the same final normalized vector defined in
\href{https://github.com/openai/math/blob/adc7f1241b42e322a6451854ab7e4b4c146bf78a/preprints/Polynomial-PEPS-approximation-of-gapped-square-grid-ground-states-September-24-2026/build/sections/03-patches.tex\#L91-L99}{the variational problem}.

\begin{definition}[One-coordinate variation and insertion]
    \label{def:peps_patch_coordinate_insertion}
    \lean{
        TNLean.PEPS.regularizedPatchCoordinateUpdate,
        TNLean.PEPS.regularizedPatchInsertion,
        TNLean.PEPS.regularizedPatchInsertion_apply,
        TNLean.PEPS.regularizedPatchInsertion_self,
        TNLean.PEPS.regularizedPatchUnitaryCurve
    }
    \leanok
    \uses{def:peps_regularized_patch_objective}
    Fix an index $j$ and a unitary $U$ on $\mathcal H_{X_j}$.
    Define $x^{j,U}$ by replacing $x_j$ with $Ux_jU^*$ and leaving every
    other coordinate unchanged. Put
    \begin{align}
        P_j^+           & =L_{m-1}\cdots L_{j+1},            &
        P_j^-           & =L_{j-1}\cdots L_0,\notag            \\
        \mathcal I_j(K) & =P_j^+\widehat K_{X_j}P_j^-\Omega.
        \label{eq:peps_coordinate_insertion}
    \end{align}
    Empty products are identities. The map $\mathcal I_j$ is complex linear,
    and $\mathcal I_j(K_j)=\Psi(x)$.
    For $B^*=-B$ and real $t$, the exponential $U(t)=e^{tB}$ is unitary.
    Equal regions at different indices still have independent coordinates.
\end{definition}

\begin{theorem}[Feasibility and covariance]
    \label{thm:peps_patch_coordinate_covariance}
    \lean{
        TNLean.PEPS.regularizedPatchCoordinateUpdate_mem,
        TNLean.PEPS.regularizedPatchCoordinateUpdate_rpow,
        TNLean.PEPS.regularizedPatchOutput_coordinateUpdate
    }
    \leanok
    \uses{def:peps_patch_coordinate_insertion}
    For $x\in\mathcal D$, $b>0$, arbitrary real weights, and a unitary $U$
    on $\mathcal H_{X_j}$, one has $x^{j,U}\in\mathcal D$ and
    \begin{align}
        (Ux_jU^*+bI)^{-a_j/2} & =UK_jU^*,
        \label{eq:peps_coordinate_covariance_rpow}      \\
        \Psi(x^{j,U})         & =\mathcal I_j(UK_jU^*).
        \label{eq:peps_coordinate_covariance_apply}
    \end{align}
\end{theorem}
\begin{proof}
    \leanok
    Unitary conjugation preserves positivity and trace. It fixes the identity,
    so $Ux_jU^*+bI=U(x_j+bI)U^*$. Unitary covariance of the real power gives
    (\ref{eq:peps_coordinate_covariance_rpow}). Every factor except the one at
    index $j$ is unchanged, giving~(\ref{eq:peps_coordinate_covariance_apply}).
    The argument also applies to singular $x_j$ because $x_j+bI$ is positive
    definite on the full regional space.
\end{proof}

\begin{theorem}[Ordered first variation]
    \label{thm:peps_patch_first_variation}
    \lean{
        TNLean.PEPS.hasDerivAt_regularizedPatchOutput_coordinateUpdate,
        TNLean.PEPS.regularizedPatchFirstVariation_eq_zero
    }
    \leanok
    \uses{def:peps_patch_coordinate_insertion}
    Let $x\in\mathcal D$, $b>0$, and $B^*=-B$ on $\mathcal H_{X_j}$.
    For arbitrary real weights, the actual filtered output satisfies
    \begin{align}
        \left.\frac{d}{dt}\Psi(x^{j,e^{tB}})\right|_{t=0}
         & =\mathcal I_j(BK_j-K_jB).
        \label{eq:peps_coordinate_derivative}
    \end{align}
    If $F(x)\le F(y)$ for every $y\in\mathcal D$, then
    \begin{align}
        \operatorname{Re}\langle\Psi(x),\mathcal I_j(BK_j-K_jB)\rangle=0.
        \label{eq:peps_coordinate_first_variation}
    \end{align}

    Set $D_j=\widehat{BK_j-K_jB}_{X_j}$ and
    $\Delta_j=P_j^+D_jP_j^-\Omega$. Every line in the following diagram
    carries the full space $\mathcal H_V$.
    % Formula: (Delta_j)_sigma = sum_{alpha,beta,gamma}
    % (P_j^+)_{sigma,alpha} (D_j)_{alpha,beta}
    % (P_j^-)_{beta,gamma} Omega_gamma.
    % Source: 03-patches.tex lines 134--149, adc7f1241b42e322a6451854ab7e4b4c146bf78a.
    % Ink-to-index map: every wire is a global configuration index.
    % D_j is the identity extension of the regional commutator; P_j^+ and
    % P_j^- are global ordered products, not operators on X_j.
    % Contractions: alpha, beta, gamma, in that left-to-right order.
    % Open leg: sigma on the left. Boundary signature: phys:w.
    % Action begins at Omega on the right: P_j^-, then D_j, then P_j^+.
    % No tensor-factor decomposition, nesting, or commutation is assumed.
    \begin{center}
        \tenkzkernel
        \begin{tenkz}[rows={wire}, cols=7, bonds=none]
            \tn[at=(1,1), name=coordinateouter, skin=box,
                ports={180:physical, 0:physical}]{P_j^+}
            \tn[at=(1,3), name=coordinateinsert, skin=box,
                ports={180:physical, 0:physical}]{D_j}
            \tn[at=(1,5), name=coordinateinner, skin=box,
                ports={180:physical, 0:physical}]{P_j^-}
            \tn[at=(1,7), name=coordinatestate, skin=box,
                ports={180:physical}]{\Omega}
            \tnwire[name=coordinateoutput]{open w}{coordinateouter.180}
            \tnwire[name=coordinatealpha]{coordinateouter.0}{coordinateinsert.180}
            \tnwire[name=coordinatebeta]{coordinateinsert.0}{coordinateinner.180}
            \tnwire[name=coordinategamma]{coordinateinner.0}{coordinatestate.180}
            \tnmark[form=label, label pos=n]{on coordinateoutput 0.5}{$\sigma$}
            \tnmark[form=label, label pos=n]{on coordinatealpha 0.5}{$\alpha$}
            \tnmark[form=label, label pos=n]{on coordinatebeta 0.5}{$\beta$}
            \tnmark[form=label, label pos=n]{on coordinategamma 0.5}{$\gamma$}
        \end{tenkz}
    \end{center}
    Starting at $\Omega$ on the right, $P_j^-$ acts first, $D_j$ second,
    and $P_j^+$ last. The two outer products need not be supported on $X_j$.
\end{theorem}
\begin{proof}
    \leanok
    \uses{thm:peps_patch_coordinate_covariance}
    Differentiate $e^{tB}K_je^{-tB}$ at zero to obtain $BK_j-K_jB$.
    Apply the continuous linear map $\mathcal I_j$ to obtain
    (\ref{eq:peps_coordinate_derivative}). The entire real curve is feasible.
    Thus $t\mapsto\|\Psi(x^{j,e^{tB}})\|^2$ has a local minimum at zero,
    and its derivative there is
    $2\operatorname{Re}\langle\Psi(x),\mathcal I_j(BK_j-K_jB)\rangle$.
    Fermat's theorem gives~(\ref{eq:peps_coordinate_first_variation}).
\end{proof}

\begin{definition}[Canonical global coordinates and marginals]
    \label{def:peps_patch_canonical_marginal}
    \lean{
        TNLean.PEPS.dependentGlobalConfigEquiv,
        TNLean.PEPS.dependentGlobalConfigIsometry,
        TNLean.PEPS.dependentGlobalConfigIsometry_apply,
        TNLean.PEPS.normalizedRegularizedPatchMarginal
    }
    \leanok
    \uses{def:peps_regularized_patch_minimum}
    Identify configurations indexed by the full region $V$ with global
    configurations by forgetting the redundant membership condition. This
    bijection induces a unitary coordinate identification $J$.
    For the same final vector $\phi=\widehat\Psi(x)$ and every region
    $X\subseteq V$, define
    \begin{align}
        \rho_X=\operatorname{Tr}_{V\setminus X}
        \ketbra{J\phi}{J\phi}.
        \label{eq:peps_canonical_marginal}
    \end{align}
    At a feasible minimizer with nonnegative weights,
    $\phi=N^{-1}M(x)\Omega$. No intermediate filtered vector enters
    (\ref{eq:peps_canonical_marginal}).
\end{definition}

\begin{theorem}[Regional expectation and density properties]
    \label{thm:peps_patch_marginal_expectation}
    \lean{
        TNLean.PEPS.inner_dependentRegionOperatorLift_eq_trace_reducedPure,
        TNLean.PEPS.normalizedRegularizedPatchMarginal_posSemidef,
        TNLean.PEPS.trace_normalizedRegularizedPatchMarginal,
        TNLean.PEPS.inner_normalizedRegularizedPatchOutput_lift_eq_trace
    }
    \leanok
    \uses{def:peps_patch_canonical_marginal}
    For every global vector $\xi$, region $X$, and arbitrary complex
    operator $K$ on $\mathcal H_X$,
    \begin{align}
        \langle\xi,\widehat K_X\xi\rangle
         & =\tr\left(\operatorname{Tr}_{V\setminus X}
        \ketbra{J\xi}{J\xi}\,K\right).
        \label{eq:peps_complex_expectation}
    \end{align}
    In particular, $\langle\phi,\widehat K_X\phi\rangle=\tr(\rho_XK)$.
    Each $\rho_X$ is positive semidefinite. If $a_j\ge0$, $b>0$,
    $\|\Omega\|=1$, and $x\in\mathcal D$, then $\tr\rho_X=1$.
    These assertions include the empty region.
\end{theorem}
\begin{proof}
    \leanok
    \uses{thm:peps_regularized_patch_minimum}
    In regional coordinates, $\widehat K_X=K\otimes I$ and the reduced
    matrix is the partial trace of $\ketbra{J\xi}{J\xi}$ over the second
    factor. The partial-trace identity gives
    $\tr(\operatorname{Tr}_{V\setminus X}(P)K)=\tr(P(K\otimes I))$.
    For $P=\ketbra{J\xi}{J\xi}$ this is
    (\ref{eq:peps_complex_expectation}); no real part or Hermiticity
    assumption is needed. Partial trace preserves positivity and trace.
    Finally $J$ preserves the norm and $\|\phi\|=1$, so $\tr\rho_X=1$.
\end{proof}
